\section{Overview: Smoothness Classes and Covariate Adjustment}\label{sec:overview}

This section presents the central theoretical argument of the paper.
The choice of smoothness class determines whether covariate adjustment can reduce worst-case bias in regression discontinuity designs.
I establish a general negative result under isotropic (pooled) smoothness and contrast it with the positive results obtainable under anisotropic (separated) smoothness.
The discrete groups setting and the continuous covariate setting are then developed as two applications of this principle.


\subsection{The General Principle}\label{sec:overview:principle}

Consider an RDD with running variable $X$, outcome $Y$, and a predetermined covariate $Z$.
The regression function $m(x,z) = \E[Y \mid X = x, Z = z]$ is bivariate; the pooled regression $\bar{m}(x) = \E[Y \mid X = x]$ integrates over $Z$.

\paragraph{Target parameter.}
Following \textcite{imbens_optimized_2019}, I consider a convex weighted average treatment effect as the target parameter:
\begin{equation*}
	\tau_w \defeq \int \tau(c, z)\, w(z)\, dz,
\end{equation*}
where $w(z)$ is a researcher-chosen weight function (e.g., uniform over the support of $Z$).
This choice has two key advantages.
First, it places the analysis directly in the framework of optimized RD designs, where \textcite{imbens_optimized_2019} recommend estimating convex weighted average treatment effects using the minimax estimator.
Second, it eliminates the need to estimate the conditional density $f_{Z|X}(z|c)$: the covariate-adjusted estimator averages segment-specific estimates $\hat\tau_j$ with \emph{known} weights $w_j$, requiring no density estimation.
By contrast, the base RD estimator implicitly targets $\bar{m}(x) = \int m(x,z) f_{Z|X}(z|x)\, dz$, whose bias inherently depends on the density through the compositional channel.

The question is: \emph{under what conditions on the smoothness of $m(x,z)$ can covariate adjustment reduce worst-case bias?}


\subsection{Negative Result: Isotropic / Pooled Smoothness}\label{sec:overview:negative}

Under an isotropic smoothness class, a single bound constrains the joint variation of $m(x,z)$:
\begin{equation*}
	\norm{\nabla m}_2 \leq M, \qquad \text{i.e.,} \quad \left(\frac{\partial m}{\partial x}\right)^2 + \left(\frac{\partial m}{\partial z}\right)^2 \leq M^2.
\end{equation*}
Equivalently, in the pooled formulation used by standard RD inference (e.g., RDHonest), a single constant $M_P$ bounds the derivative of the pooled regression:
$\abs{\bar{m}'(x)} \leq M_P$.

Under either formulation, the curvature budget is \emph{fungible} across directions.
Any estimator that controls for $Z$ --- thereby becoming insensitive to $z$-variation --- simply prompts the worst-case function to reallocate its entire curvature budget to the $x$-direction.
The worst-case function sets $\partial m / \partial z = 0$ and $\partial m / \partial x = M$, which does not vary in $z$ at all.
Against such a function, the $Z$-information is useless, and the worst-case bias equals that of the base RD estimator.

\begin{theorem}[Informal: Covariate Adjustment is Pointless under Isotropic Smoothness]\label{thm:negative_informal}
	Under the isotropic Lipschitz class $\mathcal{F}_1^{\mathrm{iso}}(M)$, for any linear estimator $\hat\tau$ that may depend on $(Y_i, X_i, Z_i)$:
	\begin{equation*}
		\wcb(\hat\tau, \mathcal{F}_1^{\mathrm{iso}}(M)) \geq \wcb(\hat\tau^*_{\mathrm{base}}, \mathcal{F}_1^{\mathrm{iso}}(M)),
	\end{equation*}
	where $\hat\tau^*_{\mathrm{base}}$ is the minimax optimal base RD estimator that ignores $Z$.
	No covariate-adjusted estimator can achieve lower worst-case bias.
\end{theorem}

The proof follows from constructing the worst-case function: for any estimator $\hat\tau$, the adversary chooses $f^*(x,z) = M \cdot x$ (or the appropriate worst-case profile in $x$ alone), which is constant in $z$.
Against this function, $Z$ carries no information, so all $Z$-adjustment is wasted.


\subsection{Positive Result: Anisotropic / Separated Smoothness}\label{sec:overview:positive}

Under the anisotropic class, smoothness in each direction is bounded separately:
\begin{equation*}
	\abs*{\frac{\partial m}{\partial x}} \leq M_x, \qquad \abs*{\frac{\partial m}{\partial z}} \leq M_z.
\end{equation*}
The $M_z$ budget \emph{cannot} be redirected into the $x$-direction.
When covariate adjustment neutralizes the $z$-channel of the bias, the $M_z$ budget is stranded: it can only generate residual bias proportional to the resolution of the adjustment (segment width $\delta$ or kernel bandwidth $g$), which is small.
The worst-case bias of the adjusted estimator therefore depends primarily on $M_x$, yielding a strict improvement over the base RD whenever $M_z > 0$ and the compositional channel is active.

This is the right framework for studying covariate adjustment: the smoothness of $m$ in $x$ is a property of the data-generating process that should not respond to the researcher's estimation strategy.
The anisotropic class respects this; the isotropic class does not.


\subsection{Coupled vs.\ Separated Smoothness: A Unifying Perspective}\label{sec:overview:epistemology}

The distinction between coupled and separated smoothness classes is the central conceptual contribution of the paper.
It applies both to the bivariate function $m(x,z)$ (2D setting) and to the pooled regression $\bar{m}(x)$ (1D setting).
The same epistemological argument governs both.

\begin{center}
\begin{tabular}{lcc}
\toprule
& \textbf{Coupled} & \textbf{Separated} \\
\midrule
\textbf{2D} (bivariate $m(x,z)$) & Isotropic: $\norm{\nabla m}_2 \leq M$ & Anisotropic: $\abs{\partial m / \partial x} \leq M_x,\; \abs{\partial m / \partial z} \leq M_z$ \\
\textbf{1D} (pooled $\bar{m}(x)$) & Pooled: $\abs{\bar{m}'(x)} \leq M_P$ & Component-wise: $M_x, M_z, L$ separately \\
\bottomrule
\end{tabular}
\end{center}

In both rows, the coupled class bundles distinct sources of variation into a single bound.
The separated class keeps them apart, preserving the directional structure that covariate adjustment exploits.
The coupled class arises from the separated class by collapsing --- via the gradient norm in 2D, or via the triangle inequality and integration in 1D --- and this collapse is lossy.


\subsubsection*{2D: Isotropic vs.\ Anisotropic}

\paragraph{Two types of ignorance.}
The standard justification for the isotropic class in the nonparametric literature is that the researcher ``does not know which direction is smoother.''
This justification conflates two distinct types of ignorance:
\begin{enumerate}
	\item \textbf{Unknown directional magnitudes:} the researcher does not know whether $M_x > M_z$ or $M_x < M_z$.
	\item \textbf{Unknown coupling:} the researcher does not know whether a large $\partial m / \partial x$ constrains the feasible values of $\partial m / \partial z$.
\end{enumerate}
The first type of ignorance --- not knowing which direction is smoother --- justifies setting equal bounds: $M_x = M_z = M$.
This gives the anisotropic class $\mathcal{F}_1^{\mathrm{aniso}}(M, M)$, whose constraint set in gradient space is a \emph{square}.
The square is symmetric across directions: it treats both coordinates identically.
Crucially, it does \emph{not} couple the partial derivatives --- each bound holds independently.

The isotropic class $\mathcal{F}_1^{\mathrm{iso}}(M)$ imposes more than directional symmetry: it constrains the gradient to lie in a \emph{disk}, which is strictly contained in the square.
The disk couples the partial derivatives --- $(\partial m / \partial x)^2 + (\partial m / \partial z)^2 \leq M^2$ --- so that a large value of one partial forces a small value of the other.
This coupling is an additional assumption beyond symmetry, and it is this coupling that makes the curvature budget fungible and renders covariate adjustment pointless (\cref{sec:overview:negative}).

The distinction matters: a researcher who is ignorant about directional magnitudes but believes the partial derivatives are independent should use the square, not the disk.
Only a researcher who additionally believes the partials \emph{may be coupled} --- and is completely agnostic about the form of that coupling --- needs the disk.

\paragraph{The isotropic class as an envelope of coupled-partial classes.}
To see why, consider the coupling question explicitly.
Any specific coupling model defines a constraint set in gradient space: a rectangle (no coupling, anisotropic), a diamond ($\ell_1$ coupling: $\abs{\partial m / \partial x} + \abs{\partial m / \partial z} \leq C$), or some other convex shape.
If the researcher \emph{knows} the coupling structure, the appropriate smoothness class is the corresponding coupled-partial class.
For any fixed coupling, this class is strictly contained in the isotropic class --- and therefore yields tighter inference (\cref{prop:iso_envelope}).

The isotropic class arises when the researcher is \emph{agnostic} about the coupling:
\begin{enumerate}
	\item the researcher believes there may or may not be coupling between the partial derivatives,
	\item the researcher does not know the form or strength of any coupling, and
	\item this uncertainty can vary locally across the domain --- at each point, the coupling structure may differ.
\end{enumerate}
Under all three conditions, the feasible gradient set at each point is the envelope of all possible coupling structures consistent with the researcher's bounds.
This envelope is a disk of radius $\sqrt{M_x^2 + M_z^2}$, recovering the isotropic class $\mathcal{F}_1^{\mathrm{iso}}(\sqrt{M_x^2 + M_z^2})$.
Because the isotropic class contains every coupled-partial class (including the anisotropic class as the special case of no coupling), it represents maximal ignorance about the coupling structure.

The cost of this agnosticism is the inflation from $M$ to $M\sqrt{2}$ (in the symmetric case): the disk of radius $M\sqrt{2}$ circumscribes the square of side $2M$.
This inflation gives the adversary strictly more room, increasing worst-case bias and minimax risk.
Any information about the coupling --- including the information that there is no coupling --- allows the researcher to shrink back from the inflated disk to a tighter class and recover the lost precision.

\paragraph{Any resolved coupling beats the isotropic class.}
The following proposition formalizes this claim.

\begin{definition}[Gradient-constrained smoothness class]\label{def:gradient_class}
	Let $\mathcal{C} \subset \mathbb{R}^2$ be a compact, convex set containing the origin.
	The gradient-constrained smoothness class induced by $\mathcal{C}$ is
	\begin{equation*}
		\mathcal{F}(\mathcal{C}) \defeq \bigl\{ f : \mathbb{R}^2 \to \mathbb{R} \;\big|\; \nabla f(x,z) \in \mathcal{C} \;\; \text{for all } (x,z) \bigr\}.
	\end{equation*}
\end{definition}

This definition nests all classes considered above.
The isotropic class corresponds to $\mathcal{C} = B_2(R)$ (the disk of radius $R$), the anisotropic class to $\mathcal{C} = [-M_x, M_x] \times [-M_z, M_z]$ (a rectangle), and any coupling model to the appropriate convex shape.

\begin{proposition}[The isotropic class is the unique coupling-invariant class]\label{prop:iso_envelope}
	Let $R > 0$ and let $B_2(R) = \{ v \in \mathbb{R}^2 : \norm{v}_2 \leq R \}$.
	\begin{enumerate}
		\item \textbf{(Envelope.)} The isotropic class is the envelope of all coupled-partial classes with bounded gradient magnitude:
		\begin{equation*}
			\mathcal{F}(B_2(R)) = \bigcup \bigl\{ \mathcal{F}(\mathcal{C}) : \mathcal{C} \subseteq B_2(R),\; \mathcal{C} \text{ compact, convex, } 0 \in \mathcal{C} \bigr\}.
		\end{equation*}
		\item \textbf{(Strict improvement.)} For any compact convex $\mathcal{C} \subsetneq B_2(R)$ with $0 \in \mathcal{C}$ and nonempty interior:
		\begin{equation*}
			\min_{\hat\tau \in \mathcal{L}} \sup_{f \in \mathcal{F}(\mathcal{C})} \mathrm{MSE}(\hat\tau, f) \;<\; \min_{\hat\tau \in \mathcal{L}} \sup_{f \in \mathcal{F}(B_2(R))} \mathrm{MSE}(\hat\tau, f).
		\end{equation*}
	\end{enumerate}
\end{proposition}

\begin{proof}[Proof sketch]
	\emph{Part 1.}
	The inclusion $\supseteq$ is immediate: for any $\mathcal{C} \subseteq B_2(R)$, $\mathcal{F}(\mathcal{C}) \subseteq \mathcal{F}(B_2(R))$.
	For $\subseteq$: take any $f \in \mathcal{F}(B_2(R))$.
	At each point $(x,z)$, $\nabla f(x,z)$ is a single vector in $B_2(R)$.
	Set $\mathcal{C} = \{ \nabla f(x,z) : (x,z) \in \mathbb{R}^2 \}$, the image of the gradient map.
	Its closed convex hull is a compact convex subset of $B_2(R)$, so $f \in \mathcal{F}(\overline{\mathrm{conv}}(\mathcal{C}))$, which appears in the union.

	\emph{Part 2.}
	The weak inequality $\leq$ follows from $\mathcal{F}(\mathcal{C}) \subsetneq \mathcal{F}(B_2(R))$ and monotonicity of minimax risk in the function class.
	For strict inequality, it suffices to show that the worst-case bias is strictly lower over $\mathcal{F}(\mathcal{C})$ than over $\mathcal{F}(B_2(R))$.

	Since $\mathcal{C} \subsetneq B_2(R)$ is compact and convex with nonempty interior, there exists a direction $\vec{v} \in \mathbb{R}^2$ with $\norm{\vec{v}}_2 = 1$ such that
	\begin{equation*}
		\sup_{g \in \mathcal{C}} \langle g, \vec{v} \rangle = h_{\mathcal{C}}(\vec{v}) < R = h_{B_2(R)}(\vec{v}),
	\end{equation*}
	where $h_{\mathcal{C}}$ denotes the support function of $\mathcal{C}$.
	The support function gives the maximal directional derivative: for any $f \in \mathcal{F}(\mathcal{C})$, the directional derivative along $\vec{v}$ satisfies $D_{\vec{v}} f \leq h_{\mathcal{C}}(\vec{v}) < R$.

	The worst-case bias of a linear estimator $\hat\tau$ over $\mathcal{F}(\mathcal{C})$ depends on the worst-case function, which maximizes the bias by concentrating its gradient in the direction that the estimator is most sensitive to.
	Under $\mathcal{F}(B_2(R))$, the adversary can always achieve directional derivative $R$ in any direction.
	Under $\mathcal{F}(\mathcal{C})$, at least one direction is strictly constrained: $h_{\mathcal{C}}(\vec{v}) < R$.
	The minimax estimator over $\mathcal{F}(\mathcal{C})$ can exploit this directional heterogeneity --- either by orienting its weight structure to be most sensitive in the less-constrained directions, or by using information (e.g., segmentation) to neutralize the constrained direction.
	In either case, the minimax estimator over $\mathcal{F}(\mathcal{C})$ achieves a strictly lower worst-case MSE than the minimax estimator over $\mathcal{F}(B_2(R))$, which cannot exploit any directional structure because the disk is rotationally symmetric.
\end{proof}

The proposition says: \emph{the isotropic class is the unique smoothness class that carries no information about the coupling structure.}
Any departure from complete agnosticism --- in any direction, by any amount --- strictly reduces minimax risk and enables tighter inference.
The anisotropic class, the $\ell_1$ class, and any other coupled-partial class are all special cases.

\paragraph{Scope beyond RDD.}
\Cref{def:gradient_class} and \cref{prop:iso_envelope} are stated for general gradient-constrained smoothness classes and make no reference to the RDD setting.
Isotropic smoothness classes are the default throughout nonparametric estimation: the foundational minimax results \parencite{stone_optimal_1980}, standard textbook treatments \parencite{tsybakov_introduction_2009}, bias-aware inference \parencite{armstrong_simple_2020}, and virtually all applied nonparametric work in econometrics use isotropic Hölder or Sobolev classes.
In each of these settings, the isotropic class silently imports the coupling assumption identified above --- and the ``two types of ignorance'' distinction (\cref{sec:overview:epistemology}) applies with full force.

The proposition implies that for \emph{any} nonparametric estimation problem --- density estimation, regression function estimation, functional estimation --- replacing the isotropic class with a class that reflects known (or assumed) independence of the partial derivatives yields a strict minimax improvement.
The improvement is quantifiable: in the symmetric case, the inflation from $M$ to $M\sqrt{2}$ translates directly into wider confidence intervals and higher worst-case MSE.
This is not a small correction: the $\sqrt{2}$ factor enters the leading bias term, and the MSE scales with the square of the bias.

The nonparametric literature has studied specific smoothness classes --- isotropic Hölder/Sobolev and anisotropic Nikol'skii/Besov (\cref{sec:appendix:smoothness_classes}) --- but does not unify them into a single framework parameterized by the constraint set $\mathcal{C}$, nor does it identify the isotropic class as the unique worst case arising from coupling agnosticism.
The proposition fills this gap: it characterizes exactly what the isotropic assumption costs and what any departure from it buys.

\paragraph{High-dimensional implications.}
The cost of the isotropic assumption grows with the dimension $d$.
In $\mathbb{R}^d$, the isotropic class constrains the gradient to lie in a $d$-dimensional ball of radius $R$; the anisotropic class constrains it to a $d$-dimensional box $[-M_1, M_1] \times \cdots \times [-M_d, M_d]$.
In the symmetric case ($M_j = M$ for all $j$), the ball of radius $M\sqrt{d}$ circumscribes the box of side $2M$.
The inflation factor is $\sqrt{d}$, which enters the leading bias term: for $d = 100$, this is a factor of $10$ in worst-case bias and potentially two orders of magnitude in worst-case MSE.
Part of what the literature calls the ``curse of dimensionality'' under isotropic smoothness is therefore self-inflicted: the adversary exploits fungibility across $d$ directions, concentrating curvature in whichever direction the estimator is most sensitive to.
Under the anisotropic class, the minimax estimator uses direction-dependent bandwidths that neutralize this reallocation, and the effective dimensionality may be lower if some directions are smoother than others.

In nonparametric estimation, the connection to anisotropic classes is already explicit for some estimators.
Wavelet estimators have been analyzed under anisotropic Besov classes, where they achieve adaptive minimax rates with direction-dependent resolution levels (Kerkyacharian, Lepski, and Picard, 2001).
Kernel estimators with product kernels use direction-dependent bandwidths, which corresponds to rectangular neighborhoods in gradient space --- the anisotropic class.
In both cases, the estimators exploit the separated structure of the anisotropic class to avoid the fungibility cost of the isotropic class.

\paragraph{The anisotropic class as the natural default.}
The anisotropic class corresponds to no coupling: the partial derivatives are independently bounded.
This is the natural default when the researcher has no reason to believe that smoothness in $x$ constrains smoothness in $z$.
In the RDD context, the running variable $X$ and the covariate $Z$ play fundamentally different roles --- the smoothness of $m$ in $x$ (the treatment assignment direction) has no obvious connection to its smoothness in $z$ (a predetermined characteristic).
Absent a specific economic model that predicts coupling, independent bounds are the appropriate assumption.
The isotropic class is justified only under complete agnosticism about coupling --- and even then, any resolution of that agnosticism leads to a class that enables tighter inference.


\subsubsection*{1D: Pooled vs.\ Component-wise}

The same logic applies when the bivariate function $m(x,z)$ is collapsed into the pooled regression $\bar{m}(x) = \int m(x,z) f_{Z|X}(z|x)\, dz$ by integrating over $Z$.

\paragraph{The pooled class.}
Standard RD inference (e.g., RDHonest, Armstrong and Kolesár 2020) places a single bound on the pooled derivative: $\abs{\bar{m}'(x)} \leq M_P$.
This bundles three distinct sources of variation into one number:
\begin{equation*}
	\bar{m}'(x) = \underbrace{\int \frac{\partial m}{\partial x}(x,z)\, f(z|x)\, dz}_{\text{smoothing (bounded by } M_x\text{)}} + \underbrace{\int m(x,z)\, \frac{\partial f}{\partial x}(z|x)\, dz}_{\text{compositional (depends on } M_z, L\text{)}}.
\end{equation*}
The bound $M_P$ merges the smoothing contribution ($M_x$), the covariate variation ($M_z$), and the density variation ($L$) into a single opaque constant.
The curvature budget is fungible: the adversary can allocate $M_P$ across sources however it wants, and the researcher cannot distinguish what fraction of $M_P$ is ``real'' smoothness vs.\ compositional variation vs.\ density effects.

\paragraph{The component-wise class.}
The separated approach states each belief directly:
\begin{equation*}
	\mathcal{F}^{\mathrm{cw}}(M_x, M_z, L) \defeq \Bigl\{ (m, f_{Z|X}) \;\Big|\; \abs{\partial m / \partial x} \leq M_x,\; \abs{\partial m / \partial z} \leq M_z,\; f_{Z|X} \text{ satisfies smoothness bound } L \Bigr\}.
\end{equation*}
The pooled bound $M_P$ is a \emph{derived} quantity, obtained via the triangle inequality:
\begin{equation*}
	M_P \leq M_x + \text{(compositional term depending on $M_z$ and $L$)}.
\end{equation*}
This parallels $M_P = M_L + C_{\Delta m} L$ in the discrete groups setting.

\paragraph{Epistemological parallel.}
The pooled class relates to the component-wise class exactly as the isotropic class relates to the anisotropic class:
\begin{itemize}
	\item The pooled class arises from the component-wise class by collapsing --- applying the triangle inequality to merge the separate bounds into $M_P$.
	This collapse is lossy: it introduces fungibility across sources.
	\item The pooled class is justified only if the researcher genuinely cannot distinguish the separate contributions --- i.e., if the allocation of $M_P$ across smoothing, compositional, and density channels is itself uncertain.
	\item In practice, the researcher typically has separate knowledge about each source: the smoothness of outcomes ($M_x, M_z$) is distinct from the variation of the covariate distribution ($L$).
	Collapsing them into $M_P$ discards this structure for no substantive reason.
\end{itemize}

The component-wise class is therefore the \emph{primitive}: it directly represents the researcher's beliefs about each component.
The pooled class is a derived, coarsened version that loses information.
This is why the ``common parameter space'' in the groups paper is not an ad hoc construction --- it is the natural starting point from which the pooled bound $M_P$ is derived.


\paragraph{Asymmetric role of the density constraint.}
The density constraint $L$ enters the component-wise class only because the base RD needs it: its bias depends on $\bar{m}(x)$, whose derivative involves the compositional term.
Without a bound on how the density varies with $x$, the base RD's worst-case bias is unbounded.

The covariate-adjusted estimator, by contrast, does not need the density constraint at all.
Because the target parameter $\tau_w = \int \tau(c,z)\, w(z)\, dz$ uses researcher-chosen weights $w(z)$, the segmented estimator averages $\hat\tau_j$ with known weights --- no density estimation required.
Its worst-case bias depends only on the anisotropic primitives:
\begin{equation*}
	\wcb(\hat\tau_{\mathrm{seg}}) \sim M_x \cdot h + M_z \cdot \delta.
\end{equation*}
The density plays no role.
This contrasts sharply with the base RD, whose bias depends on the inflated $M_P(M_x, M_z, L)$.
The comparison is therefore asymmetric: the density constraint is a burden only for the pooled approach.

This also eliminates the weight estimation bias term that appears in the groups paper (Proposition~3) from estimating $\hat{g}(c) = f_{Z|X}(1|c)$.
With known weights, that entire source of bias vanishes.


\subsection{Main Theoretical Results}\label{sec:overview:results}

The paper establishes three theoretical results of increasing depth.

\paragraph{Result 1: Function class inclusion.}
The isotropic and anisotropic classes are related by strict nesting (\cref{sec:appendix:isotropic}):
\begin{equation*}
	\mathcal{F}_1^{\mathrm{iso}}(M) \subsetneq \mathcal{F}_1^{\mathrm{aniso}}(M, M) \subsetneq \mathcal{F}_1^{\mathrm{iso}}(M\sqrt{2}).
\end{equation*}
This is mechanical but foundational: it establishes that the anisotropic class is a strict refinement of the isotropic class --- neither too large nor too small --- and that results under the anisotropic class remain valid for any DGP satisfying isotropic smoothness.

\paragraph{Result 2: Strict minimax risk ordering.}
The minimax risk over the anisotropic class is strictly lower than over the corresponding isotropic class:
\begin{equation*}
	\min_{\hat\tau \in \mathcal{L}} \sup_{f \in \mathcal{F}_{\mathrm{aniso}}(M, M)} \mathrm{MSE}(\hat\tau, f) \;<\; \min_{\hat\tau \in \mathcal{L}} \sup_{f \in \mathcal{F}_{\mathrm{iso}}(M\sqrt{2})} \mathrm{MSE}(\hat\tau, f).
\end{equation*}
The weak inequality follows from Result~1 (monotonicity of minimax risk in the function class).
The strict inequality follows because the anisotropic class has directional structure that the minimax estimator can exploit: segmentation (or a similar mechanism) reduces worst-case bias by stranding the $M_z$ budget, which is impossible under the isotropic class due to fungibility (\cref{sec:overview:negative}).

\paragraph{Result 3: Quantifying the minimax risk reduction.}
The MSE reduction from using the anisotropic class can be characterized explicitly.
\begin{itemize}
	\item \emph{Upper bound on aniso minimax risk.}
	The segmented estimator is a feasible linear estimator over $\mathcal{F}_{\mathrm{aniso}}(M_x, M_z)$.
	Its worst-case MSE provides an upper bound on the minimax risk:
	\begin{equation*}
		\min_{\hat\tau \in \mathcal{L}} \sup_{f \in \mathcal{F}_{\mathrm{aniso}}} \mathrm{MSE}(\hat\tau, f) \;\leq\; \sup_{f \in \mathcal{F}_{\mathrm{aniso}}} \mathrm{MSE}(\hat\tau_{\mathrm{seg}}, f).
	\end{equation*}
	The right-hand side is computable: the worst-case bias is $\wcb(\hat\tau_{\mathrm{seg}}) \sim M_x \cdot h + M_z \cdot \delta$ and the variance follows from within-segment estimation with independent segments.

	\item \emph{Lower bound on iso minimax risk.}
	The minimax risk over $\mathcal{F}_{\mathrm{iso}}(M\sqrt{2})$ is characterized by Armstrong and Kolesár (2020) and Imbens and Wager (2019).
	The worst-case bias scales with $M\sqrt{2}$ in the leading term.

	\item \emph{The gap.}
	With $M_x = M_z = M$, the isotropic worst-case bias scales with $M\sqrt{2}$ while the segmented estimator's bias scales with $M$ in the dominant ($x$-direction) term.
	This gives a factor of $\sqrt{2}$ reduction in the leading bias, translating to a quantifiable MSE improvement.

	More generally, for $M_x \neq M_z$, the reduction depends on the ratio $M_z / M_x$:
	\begin{equation*}
		\text{MSE reduction} = \Phi(M_z / M_x),
	\end{equation*}
	where $\Phi$ is increasing in $M_z / M_x$.
	When $M_z / M_x$ is large (strong $z$-variation relative to $x$-variation), the compositional channel dominates and the gain from segmentation is largest.
	When $M_z = 0$ (no covariate variation), there is no compositional bias and the two classes coincide: $\Phi(0) = 0$.
\end{itemize}

This triple --- inclusion, ordering, quantification --- provides a complete characterization of what the anisotropic framework buys relative to the isotropic framework for covariate adjustment in RDDs.


\subsection{Technical Contributions Beyond the Mechanical Results}\label{sec:overview:technical}

Results~1--3 are conceptually important but technically light: the nesting, ordering, and quantification follow from known minimax machinery applied to different constraint sets.
The following results require genuine technical work.

\paragraph{General negative theorem (formal).}
The informal \cref{thm:negative_informal} must be made rigorous: for \emph{any} linear estimator $\hat\tau = \sum_i w_i(X_i, Z_i) Y_i$ whose weights may depend on $Z$, the worst-case bias over $\mathcal{F}_1^{\mathrm{iso}}(M)$ is at least as large as the worst-case bias of the minimax optimal base RD estimator that ignores $Z$.
The difficulty is that the worst-case function depends on the weight structure --- one must show that the adversary can \emph{always} concentrate curvature in $x$ regardless of how cleverly the estimator exploits $Z$-information.
This requires constructing, for each weight vector $w$, a function $f^* \in \mathcal{F}_1^{\mathrm{iso}}(M)$ that is constant in $z$ and achieves the maximal bias.
This is a clean, quotable result: \emph{under isotropic smoothness, covariate adjustment is provably pointless for any linear estimator}.

\paragraph{Modulus of continuity under the anisotropic class.}
\textcite{armstrong_simple_2020} characterize minimax optimal estimators through the modulus of continuity of the target functional over the smoothness class.
Under the isotropic class, the constraint set in gradient space is a disk; under the anisotropic class, it is a rectangle.
Computing the modulus of continuity for the anisotropic class --- and showing how it differs from the isotropic case --- would connect directly to the Armstrong--Kolesár framework.
This gives the \emph{exact} minimax risk under the anisotropic class, not just the upper bound from the segmented estimator.

\paragraph{Characterizing the minimax optimal estimator under the anisotropic class.}
Solving the Imbens--Wager / Armstrong--Kolesár weight optimization problem with the rectangular constraint (rather than the disk) would reveal:
\begin{enumerate}
	\item The exact form of the minimax optimal weights --- does the optimal estimator under the anisotropic class resemble segmentation, or does it take a different form?
	\item The sharp minimax constant, not just the rate.
	\item Whether the segmented estimator has a gap relative to the minimax optimum, and if so, how large.
\end{enumerate}
This is the deepest result and could constitute a standalone contribution.
For the present paper, characterizing the minimax estimator for the special case of local constant estimation with uniform kernel (the setting used throughout) may be a feasible target.


\subsection{Two Applications}\label{sec:overview:applications}

\paragraph{Application 1: Discrete groups.}
When $Z$ is discrete, the covariate-adjusted estimator is the Stratified RD estimator, which estimates group-specific treatment effects $\tau(c, z)$ and averages them.
With a convex weighted average target, the weights are chosen by the researcher (e.g., equal weights across groups), eliminating the need to estimate $f_{Z|X}(z|c)$.
Stratification fully eliminates the compositional channel (groups are exactly separated).
The common parameter space $\mathcal{F}^R(\mathbf{M}, \mathbf{C}, \mathbf{L})$ yields worst-case AMSE dominance of the Stratified RD over the base RD.
Under treatment effect homogeneity and differing smoothness orders for $m(x,z)$ and $g(x)$, the Stratified RD achieves a better convergence rate.
The canonical CATE $\tau(c) = \sum_z \tau(c,z) f_{Z|X}(z|c)$ requires estimating $f_{Z|X}$ and introduces additional weight estimation bias --- this is left as an extension.

\paragraph{Application 2: Continuous covariates.}
When $Z$ is continuous, the covariate-adjusted estimator is the segmented estimator, which partitions the $Z$-space into intervals, runs base RD within each segment, and averages with known weights $w_j$.
No density estimation is required.
Segmentation reduces but does not fully eliminate the compositional channel (segments have positive width $\delta > 0$).
The residual within-segment compositional bias is $O(M_z \cdot \delta)$, which is small for narrow segments.
The common parameter space $\mathcal{F}^R(M_x, M_z, L)$ yields a worst-case AMSE comparison analogous to the discrete case, with the density constraint binding only for the base RD.
